% Fragmento público generado desde propietarios íntegros.
% Residencia: 52.13.3.
% Fuente: ../incoming/radion_monografia_20260827/manuscrito/sections/13_sintesis.tex:185-247
\subsubsection{Controles negativos del sistema radional}

La construcción es refutable en puntos concretos:
\begin{enumerate}
\item Si \(S_E\notin(1,2)\), el cociente APP no es uno y el cierre cambia de
celda.
\item Si la incidencia activa no produce \(N_6=90\), o si no hay dos hojas,
\(20/81\) deja de seguirse.
\item Si el verificador recibe \(180/\pi\) o \(\epsR\) antes de construir los
operandos, la generación es circular.
\item Si la resta por tríadas no estabiliza prefijos compatibles, no existe
la sección coinductiva propuesta.
\item Si \(|C^*|\ge A\), la elipse positiva deja de existir.
\item Si \(a_Sb_S\ne2\hret\), falla el teorema de área.
\item Si \(C^*=0\), deben cumplirse \(d=0\), \(\eta=0\), \(\tau=i\) y
\(j=1728\).
\item Si el área Klein del subdisco no sigue \eqref{eq:klein-area-R}, falla la
tesis de profundidad infinita.
\item Si \(F_{\mathrm{spec}}=\epsR^\circ\) se afirma sin contratérmino independiente, la
diferencia numérica \eqref{eq:chiR} lo refuta.
\item La producción de \(12:-1\) exige conservar, junto a los cardinales
\(90/120\), los doce sectores dodecafásicos, la única costura orientada y
la acción del lector sobre esos generadores. Cambiar alguna multiplicidad
debe modificar la imagen \(12S_{90}-S_{120}\). Para la cadena fundamental
sin alteración, el coeficiente de \((1-q)^{-1}\) debe resultar \(1/24\).
\item Si la función \(j\) se usa para seleccionar signo de carga, se pierde
la invariancia \eqref{eq:j-invariance} y se incurre en un error de tipo.
\end{enumerate}

\subsubsection{Realización física conjunta del radión}

La imagen final no es un círculo corregido. Es una arquitectura de
publicaciones:
\[
\begin{aligned}
&\text{APP:}&&\text{unidad, resto, acarreo y fibra};\\
&\text{TRIT:}&&\text{régimen, orientación y hojas};\\
&\text{TPK:}&&\text{fase, retorno, dodecafase y memoria};\\
&\text{elipse:}&&\text{acción y anisotropía};\\
&\text{Klein:}&&\text{profundidad interior};\\
&\text{hélice:}&&\text{biografía del retorno};\\
&\text{MD:}&&\text{reloj, energía, carga, campo y causalidad}.
\end{aligned}
\]

La ecuación del radión reúne la cara visible:
\BigRadionEquation
La elipse reúne la cara de acción:
\[
\mathcal A(1)=\hret,\qquad \mathcal A(2\pi)=\hfull.
\]
La métrica de Klein reúne la profundidad:
\[
K=-1,\qquad \operatorname{Area}_K=\infty.
\]
La holonomía reúne la memoria:
\[
g(k+9)=g(k),\qquad B_{k+9}\ne B_k.
\]

\begin{thesisbox}
El círculo es la sombra angular; la elipse es el cuerpo de acción; Klein mide
el interior; la hélice conserva la biografía. El radión es la unidad que
permite leer las cuatro afirmaciones sobre un único estado sin reducir una a
otra.
\end{thesisbox}
% Fuente: ../incoming/radion_monografia_20260827/manuscrito/sections/B_valores_y_verificacion.tex:1-105
\subsubsection{Valores canónicos, estabilidad y verificación reproducible}

\paragraph{Valores de los invariantes internos generados}

\begin{longtable}{p{.27\textwidth} p{.61\textwidth}}
\toprule
Objeto & Valor usado en el perfil coinductivo\\
\midrule
\endfirsthead
\toprule
Objeto & Valor usado en el perfil coinductivo\\
\midrule
\endhead
\(\pih\) & \(\begin{aligned}3.1415926535897932384626433832795028841\\[-2pt]971693993751058209749445923\ldots\end{aligned}\)\\
\(e_{\HMT}\) & evaluación coinductiva del lector de Euler, sin tabla externa\\
\(\alphah\) & \(\begin{aligned}0.0072973525692838009972851054723806629\\[-2pt]995641725396427091854046825\ldots\end{aligned}\)\\
\(A\) & \(\begin{aligned}7.2973525692838009972851054723806629995\\[-2pt]641725396427\ldots{}^\circ\end{aligned}\)\\
\(C^*\) & \(2.1237383389686457242619513803933607937\ldots{}^\circ\)\\
\(\eta\) & \(0.299689683921630799684926562863927\ldots\)\\
\(\hret\) & \(1.05457181691503906113369058472430\ldots\times10^{-34}U_S\)\\
\bottomrule
\end{longtable}

\paragraph{Convergencia cuantitativa por profundidad}

\begin{center}
\small
\begin{tabular}{r l l}
\toprule
Profundidad & \(\epsR^\circ\) & lectura\\
\midrule
12 & \(5.1681282186551375597\times10^{-8}\) & preasintótica\\
18 & \(5.1383017764316644535\times10^{-8}\) & 7 cifras estables\\
24 & \(5.1383016758575394560\times10^{-8}\) & 14 cifras estables\\
30 & \(5.1383016758575282072\times10^{-8}\) & 20 cifras estables\\
36 & \(5.138301675857528207168517\times10^{-8}\) & perfil largo\\
45 & \(5.13830167585752820716851646226459\times10^{-8}\) & ancho \(<4.81\times10^{-130}\)\\
\bottomrule
\end{tabular}
\end{center}

Los retornos de profundidad \(9,18,27,36,45\) conservan la fase \(8\) y
aumentan la memoria \(0,1,2,3,4\). La razón entre anchos consecutivos tiende
a
\[
3^{-54}=1.7196982334549856127610399316004871\ldots\times10^{-26}.
\]

\paragraph{Análisis de sensibilidad estructural por supresión controlada}

\begin{longtable}{p{.46\textwidth} p{.34\textwidth}}
\toprule
Operación & Salida o diferencia unitaria\\
\midrule
\endfirsthead
\toprule
Operación & Salida o diferencia unitaria\\
\midrule
\endhead
Operador canónico & \(8.96802822044562981999\times10^{-10}\)\\
Una hoja, \(\rho=90/729\) & \(-1.3727056615960678\times10^{-5}\)\\
Fase \(\theta_1=20^\circ\) & \(+5.2359877649510169\times10^{-1}\)\\
Fase \(\theta_3=80^\circ\) & \(-5.2359877470149605\times10^{-1}\)\\
Sin canal \(e\) & \(1.8065350748904653\times10^{-3}\)\\
Torsión histórica \(\Delta_4(A,C^*)\) & \(1.0517084608768816\times10^{-9}\)\\
Operador armónico \(M_4\) menos canónico & \(1.0525500222895070\times10^{-17}\)\\
Cola sexádica menos canónico & \(2.0283936357433839\times10^{-12}\)\\
\bottomrule
\end{longtable}

La separación de escalas muestra sensibilidad estructural: el cierre depende
de la fase, las dos hojas, los canales \(\pi/e\), la torsión global y el acarreo.
Ninguna ablación reproduce el valor.

\paragraph{Certificados computacionales reproducibles}

La obra usa los siguientes propietarios ejecutables:
\begin{enumerate}
\item \sourcepath{output/RADION_TWOWAY_NONADICO_20260827/verificar_emergencia_radion_por_profundidad.py};
\item \sourcepath{output/RADION_TWOWAY_NONADICO_20260827/verificar_radion_helice_critica_tpk.py};
\item \sourcepath{output/RADION_TWOWAY_NONADICO_20260827/verificar_operador_tpk_target_free_profundidad.py};
\item el verificador consolidado de esta monografía, ubicado en
\sourcepath{certificados/verificar_monografia_radion.py}.
\end{enumerate}
Las salidas esperadas incluyen
\[
\texttt{PASS\_EMERGENCIA\_RADION\_POR\_PROFUNDIDAD},
\]
\[
\texttt{PASS\_RADION\_HELICE\_CRITICA\_TPK},
\qquad
\texttt{PASS\_MONOGRAFIA\_RADION\_HMT\_MD}.
\]

\paragraph{Distinción entre precisión computacional y exactitud matemática}

Que un programa muestre cero a 120 o 200 cifras no convierte por sí solo una
identidad residual en predicción. La exactitud matemática procede de
\eqref{eq:cierre-unitario}; la precisión arbitraria permite evaluar sus
operandos y verificar la convergencia. La no circularidad procede del orden
de construcción y de la prohibición de blancos, no del número de dígitos.

El cálculo Decimal puede dejar residuos del orden \(10^{-218}\) al combinar
ramas evaluadas con precisiones distintas. Ese número es redondeo de la
implementación; no una corrección física adicional.
% Fuente: ../incoming/radion_monografia_20260827/manuscrito/sections/D_historia_fuentes.tex:1-169
\subsubsection{Genealogía conceptual y fuentes primarias}

\paragraph{Cuanto de acción de Planck}

Planck introduce \(h\) en la ley de radiación como constante que liga energía
y frecuencia. El dato histórico decisivo para esta obra es dimensional:
\(h\) es acción. El radión no modifica ese papel; explica por qué la acción de
una vuelta y la acción por unidad de fase se relacionan mediante \(2\pi\).

\paragraph{Cuantización angular de Bohr y Sommerfeld}

Bohr usa \(h/2\pi\) en la cuantización del momento angular. Sommerfeld amplía
la cuantización a integrales de acción y órbitas keplerianas. Por ello, no es
históricamente exacto decir que Dirac inventó la constante reducida. Su
contribución pertenece a una genealogía en la que la acción angular ya estaba
presente.

La novedad HMT--MD es distinta: no toma \(h/2\pi\) como un cociente externo,
sino que construye una elipse cuya vuelta tiene área \(h\) y cuyo sector de un
radián tiene área \(\hbar\).

\paragraph{Formalismo matricial de Born y Jordan y oscilador cuántico}

Born y Jordan sitúan el principio variacional, las fases y el oscilador
armónico en el corazón de la mecánica matricial; el campo electromagnético se
descompone en una colección de osciladores. La realización LC del radión
continúa esa intuición con un dato añadido: el cociente de forma
\(e^{-\eta}\) determina la impedancia relativa mientras producto, frecuencia
y acción permanecen fijos.

\paragraph{Frecuencia angular y acción en la formulación de Dirac}

En 1925, Dirac escribe las frecuencias en radianes por unidad de tiempo y usa
relaciones de la forma
\[
xy-yx=\frac{ih}{2\pi}[x,y],
\qquad
\Delta E=\frac{h}{2\pi}\omega.
\]
La ecuación \(E=\hret\omega\) une de manera explícita acción reducida y
frecuencia angular. En el radión, esa relación recibe una geometría: \(\hret\)
es el área barrida por una unidad de fase.

\paragraph{Relaciones de indeterminación de Heisenberg, Robertson y Schrödinger}

Heisenberg formula la limitación conjunta de variables canónicamente
conjugadas. Robertson generaliza la desigualdad; Schrödinger incorpora la
covarianza. La matriz \eqref{eq:covariance} pertenece más directamente a esa
prolongación Robertson--Schrödinger que al artículo de Heisenberg aislado.

Ninguno de estos autores formula la elipse genealógica completa de esta obra.
La afirmación históricamente defendible es que proporcionan el lenguaje en el
que la elipse de acción puede reconocerse como un contorno mínimo; la
construcción del par \((A,C^*)\), los focos, Klein, la hélice y el acarreo es la
aportación HMT--MD.

\paragraph{Funcional de fase de Feynman}

La formulación espacio-temporal asigna a cada camino una fase
\(e^{iS/\hbar}\). La unidad de radión hace literal la correspondencia:
incrementar la acción en \(\hbar\) incrementa la fase en un radián. La suma de
caminos reconoce, en un lenguaje integral, el mismo principio que la elipse
expresa como área.

\paragraph{Propagaciones avanzada y retardada en Wheeler--Feynman}

La teoría del absorbedor construye una interacción clásica temporalmente
simétrica mediante combinaciones de campos retardados y avanzados. La
identidad HMT \(CU(t)C^{-1}=U(-t)\) y la inversión de rutas proporcionan un
antecedente operatorio. La promoción física requiere que el mapa entre rutas
y propagadores preserve acción, orientación, frontera y forma simpléctica.

La teoría del absorbedor no es la misma afirmación que la interpretación de
Feynman y Stückelberg de las antipartículas como propagación temporalmente
invertida. Ambas pueden reconocer la bidireccionalidad; sus objetos físicos y
condiciones de frontera son diferentes.

\paragraph{Dualidad de Hodge y separación respecto de la cohomología de Hochschild}

La operación relevante en este manuscrito es la estrella de Hodge
\(\star^2=-I\) sobre dos-formas en firma lorentziana. No se ha localizado en
el corpus activo del radión un propietario de homología de Hochschild que
actúe sobre la elipse, el acarreo o el sector \(90/120\). Introducir Hochschild
por semejanza nominal degradaría el tipado. Si una futura prolongación
construye un cociclo de deformación con dominio y codominio propios, deberá
entrar como resultado nuevo, no como retroatribución.

\paragraph{Bibliografía primaria seleccionada}

\begin{enumerate}
\renewcommand{\labelenumi}{[\arabic{enumi}]}
\item\hypertarget{radionbib:Planck1901}{}
M.~Planck,
\emph{Ueber das Gesetz der Energieverteilung im Normalspectrum},
Annalen der Physik 309 (1901), 553--563.
\url{https://doi.org/10.1002/andp.19013090310}

\item\hypertarget{radionbib:Bohr1913}{}
N.~Bohr,
\emph{On the Constitution of Atoms and Molecules},
Philosophical Magazine 26 (1913), 1--25.
\url{https://doi.org/10.1080/14786441308634955}

\item\hypertarget{radionbib:Sommerfeld1916}{}
A.~Sommerfeld,
\emph{Zur Quantentheorie der Spektrallinien, I},
Annalen der Physik 356 (1916), 1--94.
\url{https://doi.org/10.1002/andp.19163561702}

\item\hypertarget{radionbib:BornJordan1925}{}
M.~Born y P.~Jordan,
\emph{Zur Quantenmechanik},
Zeitschrift für Physik 34 (1925), 858--888.
\url{https://doi.org/10.1007/BF01328531}

\item\hypertarget{radionbib:Dirac1925}{}
P.~A.~M.~Dirac,
\emph{The Fundamental Equations of Quantum Mechanics},
Proceedings of the Royal Society A 109 (1925), 642--653.
\url{https://doi.org/10.1098/rspa.1925.0150}

\item\hypertarget{radionbib:Heisenberg1927}{}
W.~Heisenberg,
\emph{Über den anschaulichen Inhalt der quantentheoretischen Kinematik und Mechanik},
Zeitschrift für Physik 43 (1927), 172--198.
\url{https://doi.org/10.1007/BF01397280}

\item\hypertarget{radionbib:Dirac1928}{}
P.~A.~M.~Dirac,
\emph{The Quantum Theory of the Electron},
Proceedings of the Royal Society A 117 (1928), 610--624.
\url{https://doi.org/10.1098/rspa.1928.0023}

\item\hypertarget{radionbib:Robertson1929}{}
H.~P.~Robertson,
\emph{The Uncertainty Principle},
Physical Review 34 (1929), 163--164.
\url{https://doi.org/10.1103/PhysRev.34.163}

\item\hypertarget{radionbib:Schrodinger1930}{}
E.~Schrödinger,
\emph{Zum Heisenbergschen Unschärfeprinzip},
Sitzungsberichte der Preussischen Akademie der Wissenschaften (1930), 296--303.

\item\hypertarget{radionbib:WheelerFeynman1945}{}
J.~A.~Wheeler y R.~P.~Feynman,
\emph{Interaction with the Absorber as the Mechanism of Radiation},
Reviews of Modern Physics 17 (1945), 157--181.
\url{https://doi.org/10.1103/RevModPhys.17.157}

\item\hypertarget{radionbib:Feynman1948}{}
R.~P.~Feynman,
\emph{Space-Time Approach to Non-Relativistic Quantum Mechanics},
Reviews of Modern Physics 20 (1948), 367--387.
\url{https://doi.org/10.1103/RevModPhys.20.367}

\item\hypertarget{radionbib:Feynman1949}{}
R.~P.~Feynman,
\emph{The Theory of Positrons},
Physical Review 76 (1949), 749--759.
\url{https://doi.org/10.1103/PhysRev.76.749}
\end{enumerate}

\begin{narrativebox}
Los fundadores contienen las piezas: acción, frecuencia angular, variables
conjugadas, oscilador, incertidumbre, fase de caminos y reversión temporal.
La aportación del radión es reunirlas en un objeto generado: una elipse de
área \(h\), un sector de área \(\hbar\), un interior de profundidad infinita,
una orientación de hoja y un acarreo exacto de compatibilidad.
\end{narrativebox}
% Fuente: ../incoming/radion_monografia_20260827/manuscrito/sections/E_conclusion.tex:1-107
\subsubsection{Síntesis final: unidad angular con genealogía de fase, acción y memoria}

El radián no desaparece. Se vuelve visible por completo.

En su uso ordinario, un radián es una razón que no necesita recordar cómo fue
producida. El arco y el radio se cancelan dimensionalmente; la unidad parece
desnuda. Esa desnudez es una virtud para la trigonometría, pero es un olvido
para una teoría que pretende explicar la estructura del continuo y el origen
de la acción. El radión recupera lo olvidado sin alterar el valor publicado.

La genealogía comienza en APP, donde una salida conserva valor y también
hoja, cociente, resto, acarreo y frontera. El TRIT introduce régimen y
orientación. El TPK hace volver la fase sin reiniciar el estado. La conexión
nonádica construye la sección coinductiva de \(\pi\), y la dodecafase publica
la bisagra de \(50^\circ\). El par angular aporta un modo común \(A\) y una
apertura orientada \(C^*\). La incidencia de seis fases activas produce
\(90\), sus dos hojas sobre el ambiente \(729\) producen \(20/81\), y la
torsión global aporta la segunda sección.

Entonces aparece \(\varepsilon_R\). No como un decimal suplicado por el
cierre, sino como la diferencia que la estructura estaba obligada a
transportar:
\[
\epsone=\frac{20}{81}\Delta_4-(S_E-1).
\]
El acarreo publica sus cifras; la profundidad certifica su límite; la carta
angular lo reconoce como
\[
5.13830167585752820716851646226459474899\ldots
\times10^{-8}\,{}^\circ.
\]
La ecuación del radián queda cerrada porque ambos lados son publicaciones del
mismo estado:
\BigRadionEquation

La historia no termina en la igualdad. El par \((A,C^*)\) construye una
elipse positiva. Su producto de semiejes fija \(2\hret\); su cociente fija la
forma. El teorema de área demuestra que una unidad de fase barre \(\hret\) y
que la vuelta encierra \(\hfull\). Aquí adquiere contenido físico la división
por \(2\pi\): \(h\) es acción de vuelta; \(\hbar\) es acción por radión.

Los focos dejan de ser adornos de una figura. Su separación reconstruye
\(C^*\) en la carta espectral, y su distancia de Klein mide una profundidad
hiperbólica finita dentro de un dominio cuya área hiperbólica total es
infinita. La misma frontera contiene una acción finita. La aparente paradoja
es la tesis central de la geometría: una medida cierra la acción; otra abre la
profundidad.

El círculo, por tanto, no es falso. Es la proyección exacta que conserva fase
y olvida anisotropía, interior y memoria. La elipse conserva acción y forma.
Klein conserva profundidad. La hélice conserva retorno y biografía. Ninguna
imagen sustituye a las demás; todas nacen del mismo estado.

La realización de Hilbert reconoce la elipse como contorno de covarianza
mínima. El oscilador LC la reconoce como intercambio de reservas eléctrica y
magnética. El lector \(90/120\) publica permeabilidad, permitividad,
impedancia y velocidad. La incidencia seis--ocho produce una pantalla causal;
Hodge convierte orientación en dualidad de campos; Lorentz convierte campo y
carga en dinámica; Cartan--Holst recibe la holonomía en geometría
gravitatoria. El orden importa: el objeto HMT selecciona; la física
convencional realiza.

La palabra \emph{radión} queda así justificada. No significa que un foco sea
una carga positiva y el otro una carga negativa. Significa que la unidad de
fase posee orientación firmada antes de la publicación eléctrica:
\[
\mathcal A(\mathfrak r_+)=+\hret,
\qquad
\mathcal A(\mathfrak r_-)=-\hret.
\]
La carga aparece cuando Mecánica Dimensional aplica el transductor. El signo
no se añade al final; se conserva desde la hoja.

También queda explicado por qué el sector \(90/120\) parecía esconder muchas
respuestas. Incidencia, extractor dodecafásico, cola espectral, núcleo de
Barbero, torres y acarreo comparten origen. Precisamente por compartir origen
pueden compararse. Precisamente por tener operaciones distintas no deben
fundirse. La tipificación no reduce la unidad del sistema: la hace inteligible.

% REMATE_LOCAL_RADION_CIERRE
\Needspace{28\baselineskip}
La conclusión más compacta de esta obra puede escribirse en cuatro líneas:
\[
\boxed{
\begin{aligned}
\text{fase:}\quad&\theta=1,\\
\text{acción:}\quad&\mathcal A=\sigma\hret,\\
\text{memoria:}\quad&H_9(n+9)=H_9(n)+(0,1),\\
\text{compatibilidad:}\quad&1=S_E-\Phi_T+\epsone.
\end{aligned}}
\]

El radión es esa unidad cuádruple. El radián es su sombra exacta en la carta
angular. La diferencia no está en el valor de \(\pi\); está en cuánta
estructura estamos dispuestos a conservar cuando decimos que una vuelta ha
cerrado.

\vspace{2\baselineskip}
\begin{center}
\begin{minipage}{.82\textwidth}
\centering\large\itshape\color{RadNavy}
El círculo publica la fase.\\
La elipse contiene la acción.\\
Klein abre la profundidad.\\
La hélice recuerda el camino.\\[.7em]
El radión es la unidad que no obliga a elegir entre ellas.
\end{minipage}
\end{center}
